\section{Assignment 1}
\label{sec:assignment1}
\begin{exercise}
Choose any possible stream of AI application in Nepal and explore the
followings:
\begin{enumerate}
  \item AI Application domain with Scope, Challenges and Opportunities
  \item Possible Risks and risk mitigation strategy
  \item Job replacement scenario and how to deal with unemployment
  \item Advantages vs limitations lists
  \item Government role and social impact

\end{enumerate}
\end{exercise}

\begin{answer}
\textbf{AI Application Domain: Agriculture in Nepal}

\begin{enumerate}
  \item \textbf{AI Application Domain with Scope, Challenges, and Opportunities}

Artificial Intelligence in agriculture holds transformative potential for Nepal, a country where a large portion of the population relies on farming. AI applications include precision farming, disease detection, weather and market prediction, and resource optimization. Tools like drones, sensors, and machine learning algorithms can help increase yield and reduce losses.

\textbf{Scope:} Precision agriculture, automated irrigation, pest/disease diagnosis, market analytics, and advisory services using AI chatbots.

\textbf{Challenges:} Lack of digital infrastructure, low AI awareness among farmers, scarcity of localized datasets, high implementation costs, and resistance to adopting new technologies.

\textbf{Opportunities:} Enhanced crop yield, climate-resilient practices, real-time advisory services, empowerment of women farmers, and smart farming systems that make agriculture more efficient and profitable.

  \item \textbf{Possible Risks and Risk Mitigation Strategy}

\textbf{Risks:} 
\begin{itemize}
  \item Data privacy breaches due to inadequate cybersecurity.
  \item Algorithmic bias from non-representative datasets.
  \item Job displacement in manual farming roles.
  \item Dependence on foreign AI tools and platforms.
  \item Environmental harm from improper AI-based decisions.
\end{itemize}

\textbf{Mitigation Strategies:}
\begin{itemize}
  \item Establish strong data governance and privacy regulations.
  \item Invest in ethical AI design and transparent algorithms.
  \item Provide cybersecurity infrastructure and training.
  \item Encourage local AI research to reduce dependency.
  \item Mandate environmental impact assessments for AI tools.
\end{itemize}

  \item \textbf{Job Replacement Scenario and How to Deal with Unemployment}

AI tools may reduce the need for manual labor through automation (e.g., AI-powered tractors or robotic harvesting). This could lead to displacement of traditional agricultural workers.

\textbf{Strategies to Address This:}
\begin{itemize}
  \item Retrain workers in digital agriculture, AI system maintenance, and data handling.
  \item Promote agri-tech entrepreneurship and startups to create new jobs.
  \item Establish vocational training centers focused on AI in agriculture.
  \item Provide government grants and microloans for tech-driven agri-businesses.
\end{itemize}

  \item \textbf{Advantages vs. Limitations}

\begin{tabular}{|p{4cm}|p{5cm}|p{5cm}|}
\hline
\textbf{Category} & \textbf{Advantages} & \textbf{Limitations} \\
\hline
Productivity & Higher yields, efficiency in resource use, automated operations & High initial investment, limited infrastructure \\
\hline
Technology & Real-time insights, disease detection, market predictions & Lack of Nepal-specific data, technical skill gap \\
\hline
Socio-economic Impact & Supports food security, empowers women, rural upliftment & May increase inequality if access is unequal \\
\hline
Sustainability & Reduced chemical use, smart resource allocation & Risk of misuse without proper training or safeguards \\
\hline
\end{tabular}

  \item \textbf{Government Role and Social Impact}

\textbf{Government Role:}
\begin{itemize}
  \item Enforce the National AI Policy 2081/2025 with clear goals for agriculture.
  \item Invest in digital infrastructure and data centers in rural areas.
  \item Offer subsidies and incentives for AI-based farming tools.
  \item Facilitate partnerships between research institutions and tech firms.
  \item Launch AI education campaigns and pilot programs for farmers.
\end{itemize}

\textbf{Social Impact:}
\begin{itemize}
  \item Improved food security and income generation for smallholders.
  \item Empowerment of marginalized communities, especially women.
  \item Potential erosion of traditional farming knowledge if not integrated.
  \item Risk of social inequality if digital divide is not addressed.
\end{itemize}
\end{enumerate}
\end{answer}
